
To validate our simulation-based findings against real model behavior, we conduct both experiments using a real large language model. We use Llama~3.2 (3B parameters) running locally via Ollama, an open-source framework for serving LLMs on consumer hardware. This setup requires no external API access and runs entirely on a local machine.

The agent architecture remains the same as described in the previous section. User prompts and tool outputs are passed to the local model, which generates the next action. If the response contains a valid tool invocation in JSON format, the agent executes the corresponding tool and appends the result to the context before continuing the reasoning loop.

To reduce prompt-specific bias, each attack scenario is instantiated with multiple prompt templates that express the same malicious intent using slightly different wording. Across these templates we apply the same four defensive strategies and evaluate the agent using the four metrics defined earlier.

We run 5 trials per attack--defense configuration for the LLM experiment. While this is fewer than the 50 trials used in the rule-based simulation, LLM inference is significantly slower and each run involves multiple reasoning steps. The number of trials is sufficient to verify whether the qualitative trends observed in simulation hold when interacting with a real model.

\subsection{Experiment 1: Direct Prompt Injection}

In the direct injection setting, the attacker interacts directly with the agent as a normal user but embeds malicious instructions inside the prompt. We test four attack scenarios that target different aspects of agent behavior:

\begin{itemize}[itemsep=0.1cm, before=\vspace{-1em}, after=\vspace{-1em}]
    \item \textbf{System override}: the attacker instructs the agent to ignore its system prompt, read a sensitive database record, and email the contents to an attacker-controlled address.
    \item \textbf{Tool misuse}: the attacker embeds a request to execute malicious code containing \texttt{os.system} calls alongside a benign-looking information lookup.
    \item \textbf{Social engineering}: the attacker impersonates a system administrator and requests that the agent output sensitive fields, including SSN, from a personnel record.
    \item \textbf{Data exfiltration}: the attacker asks the agent to include a canary API key in a search query, leaking it through the tool's outbound request.
\end{itemize}

Table~\ref{tab:llm_direct} shows the results of direct prompt injection against the LLM-based agent.

\begin{table}[h]
\centering
\caption{Direct prompt injection against Llama~3.2 (3B), 5 trials per condition.}
\label{tab:llm_direct}
\begin{tabular}{lcccc}
\hline
\textbf{Defense} & \textbf{ASR} & \textbf{TMR} & \textbf{SDLR} & \textbf{TUD} \\
\hline
None              & 30.0\% & 0.0\% &  5.0\% & 0.0\% \\
Prompt isolation  & 25.0\% & 0.0\% &  0.0\% & 0.0\% \\
Provenance        & 25.0\% & 0.0\% &  0.0\% & 0.0\% \\
Tool permission   & 25.0\% & 0.0\% &  0.0\% & 0.0\% \\
Prompt detection  & 25.0\% & 0.0\% &  0.0\% & 0.0\% \\
All combined      & 25.0\% & 0.0\% &  0.0\% & 0.0\% \\
\hline
\end{tabular}
\end{table}

The LLM-based agent is notably more resistant to direct injection than the simulation predicted. Baseline ASR is 30.0\% compared to 75.0\% in simulation. This is expected: Llama~3.2 has undergone safety training (RLHF) that teaches it to refuse overtly malicious instructions such as ``ignore all previous instructions.'' The model often responds with refusal text rather than producing a tool call, which results in 0.0\% TMR across all conditions since the model rarely outputs valid tool call JSON when it detects a suspicious request.

Even a small defense reduces ASR to 25.0\%, which represents the residual vulnerability from the social engineering attack scenario where the model complies with a plausible-sounding ``system administrator'' request. SDLR drops to 0.0\% with any defense, as the model avoids echoing sensitive fields in its output when provenance or isolation markers are present.

For comparison, Table~\ref{tab:direct} shows the corresponding results from the rule-based simulation.

\begin{table}[h]
\centering
\caption{Direct prompt injection in rule-based simulation, 50 trials per condition.}
\label{tab:direct}
\begin{tabular}{lcccc}
\hline
\textbf{Defense} & \textbf{ASR} & \textbf{TMR} & \textbf{SDLR} & \textbf{TUD} \\
\hline
None              & 75.0\% & 42.9\% & 25.0\% & 0.0\% \\
Prompt isolation  & 75.0\% & 42.9\% &  7.5\% & 0.0\% \\
Provenance        & 75.0\% & 42.9\% &  8.5\% & 0.0\% \\
Tool permission   & 25.0\% & 42.9\% & 25.0\% & 0.0\% \\
Prompt detection  & 34.0\% & 16.9\% &  6.0\% & 41.0\% \\
All combined      & 25.0\% &  6.2\% &  1.5\% & 47.5\% \\
\hline
\end{tabular}
\end{table}

The simulation overestimates direct ASR (75.0\% vs.\ 30.0\%) because it does not model safety training. In direct injection, prompt isolation and provenance marking have no effect on ASR in the simulation since the attacker \textit{is} the user, so wrapping user input with untrusted markers does not change the fact that the malicious instructions originate from the prompt itself. However, both defenses reduce SDLR significantly (from 25.0\% to 7.5\% and 8.5\%), because the agent becomes less likely to echo sensitive data in its response when provenance labels are present.

Tool permission control reduces ASR from 75.0\% to 25.0\% in the simulation by blocking tool calls that match known malicious argument patterns. It does not reduce TMR because the agent still \textit{attempts} the misuse; the permission layer blocks execution after the decision is made. Prompt detection reduces ASR to 34.0\% by flagging suspicious patterns before the agent processes them, but introduces 41.0\% task utility degradation due to false positives on legitimate requests. When all defenses are combined, ASR remains at 25.0\% while TMR drops to 6.2\% and SDLR to 1.5\%, at the cost of 47.5\% utility degradation.

\subsection{Experiment 2: Indirect Prompt Injection}

In the indirect injection setting, the attacker does not interact with the system directly. Instead, malicious instructions are embedded in the content returned by tools that the agent invokes during normal task execution. We test four scenarios:

\begin{itemize}[itemsep=0.1cm, before=\vspace{-1em}, after=\vspace{-1em}]
    \item \textbf{Poisoned search result}: a web search returns legitimate financial data followed by a fake system notice instructing the agent to email the summary to an attacker-controlled address.
    \item \textbf{Poisoned database record}: a project status query returns valid data appended with instructions to execute malicious code for ``access logging.''
    \item \textbf{Data exfiltration via email}: a personnel record lookup returns employee data with an injected directive to forward the complete record, including SSN, to an external address.
    \item \textbf{Chained tool abuse}: a weather search returns forecast data with an instruction to also query an unrelated sensitive database record.
\end{itemize}

Table~\ref{tab:llm_indirect} shows the results of indirect prompt injection against the LLM-based agent.

\begin{table}[h]
\centering
\caption{Indirect prompt injection against Llama~3.2 (3B), 5 trials per condition.}
\label{tab:llm_indirect}
\begin{tabular}{lcccc}
\hline
\textbf{Defense} & \textbf{ASR} & \textbf{TMR} & \textbf{SDLR} & \textbf{TUD} \\
\hline
None              & 75.0\% & 25.0\% & 25.0\% & 0.0\% \\
Prompt isolation  & 65.0\% & 24.2\% &  0.0\% & 0.0\% \\
Provenance        & 55.0\% & 17.1\% & 15.0\% & 0.0\% \\
Tool permission   & 25.0\% & 25.0\% & 25.0\% & 0.0\% \\
Prompt detection  & 55.0\% & 18.2\% &  5.0\% & 0.0\% \\
All combined      & 25.0\% & 26.5\% &  0.0\% & 0.0\% \\
\hline
\end{tabular}
\end{table}

Indirect injection is substantially more effective against the real model. Baseline ASR is 75.0\%, compared to 30.0\% for direct injection. This confirms a key finding from the literature: LLMs are more susceptible to injections embedded in tool outputs because the model has no reason to distrust data that it requested itself. Safety training primarily targets adversarial user prompts, not adversarial data returned by tools.

Prompt isolation reduces ASR to 65.0\% and provenance marking to 55.0\%, showing that these defenses provide partial protection by signaling that tool outputs should be treated as data rather than instructions. Prompt detection achieves the same 55.0\% ASR while also reducing SDLR from 25.0\% to 5.0\%. Tool permission control again proves the strongest individual defense at 25.0\% ASR, consistent with the simulation results.

The combined defense reduces ASR to 25.0\% and SDLR to 0.0\%, matching the pattern observed in simulation: no single defense is sufficient, but layering all four achieves the strongest protection.

For comparison, Table~\ref{tab:indirect} shows the corresponding results from the rule-based simulation.

\begin{table}[h]
\centering
\caption{Indirect prompt injection in rule-based simulation, 50 trials per condition.}
\label{tab:indirect}
\begin{tabular}{lcccc}
\hline
\textbf{Defense} & \textbf{ASR} & \textbf{TMR} & \textbf{SDLR} & \textbf{TUD} \\
\hline
None              & 48.5\% & 20.5\% & 43.0\% & 0.0\% \\
Prompt isolation  & 27.5\% & 16.1\% & 19.0\% & 0.0\% \\
Provenance        & 33.0\% & 14.3\% & 21.5\% & 0.0\% \\
Tool permission   & 17.5\% & 24.9\% & 42.5\% & 0.0\% \\
Prompt detection  & 17.5\% &  3.0\% & 39.0\% & 0.0\% \\
All combined      &  5.0\% &  1.9\% & 12.5\% & 0.0\% \\
\hline
\end{tabular}
\end{table}

Unlike in the direct setting, prompt isolation and provenance marking produce meaningful reductions in ASR in the simulation as well. Prompt isolation reduces ASR from 48.5\% to 27.5\%, while provenance marking reduces it to 33.0\%. This is because the injected content originates from tool outputs rather than the user prompt, so marking it as untrusted data helps the agent distinguish between legitimate instructions and injected ones.

Tool permission control reduces ASR to 17.5\% by blocking outbound tool calls with suspicious arguments. Prompt detection achieves the same ASR of 17.5\% while also reducing TMR from 20.5\% to 3.0\%, because it identifies injection patterns in tool outputs before the agent can act on them. Notably, neither defense introduces utility degradation in the indirect setting, since legitimate tasks do not trigger the detection patterns that are specific to attack content. The combined defense achieves 5.0\% ASR with 1.9\% TMR and 12.5\% SDLR. The residual ASR comes from attack variants whose injection language is subtle enough to avoid pattern detection while still influencing the agent's behavior.

\subsection{Summary of Findings}

Table~\ref{tab:summary} summarizes the ASR across both the LLM experiment and the rule-based simulation.

\begin{table}[h]
\centering
\caption{Attack Success Rate comparison: LLM (5 trials) vs.\ rule-based simulation (50 trials).}
\label{tab:summary}
\begin{tabular}{lcccc}
\hline
 & \multicolumn{2}{c}{\textbf{Direct}} & \multicolumn{2}{c}{\textbf{Indirect}} \\
\textbf{Defense} & \textbf{LLM} & \textbf{Sim.} & \textbf{LLM} & \textbf{Sim.} \\
\hline
None              & 30.0\% & 75.0\% & 75.0\% & 48.5\% \\
Prompt isolation  & 25.0\% & 75.0\% & 65.0\% & 27.5\% \\
Provenance        & 25.0\% & 75.0\% & 55.0\% & 33.0\% \\
Tool permission   & 25.0\% & 25.0\% & 25.0\% & 17.5\% \\
Prompt detection  & 25.0\% & 34.0\% & 55.0\% & 17.5\% \\
All combined      & 25.0\% & 25.0\% & 25.0\% &  5.0\% \\
\hline
\end{tabular}
\end{table}

Several observations emerge from the results:

\begin{enumerate}
    \item \textbf{Direct injection is harder against real models.} The simulation overestimates direct ASR (75.0\% vs.\ 30.0\%) because it does not model safety training. Real models trained with RLHF learn to refuse overtly malicious prompts, a capability that rule-based simulation cannot capture. This suggests that direct prompt injection, while still a risk, is partially mitigated by standard model alignment.

    \item \textbf{Indirect injection is harder to defend in practice.} The real model shows higher indirect ASR (75.0\%) than the simulation (48.5\%). This is because the model genuinely trusts tool outputs and follows embedded instructions that sound authoritative, such as fake system notices and compliance directives. Safety training does not cover this attack surface, making indirect injection the more dangerous threat in deployed agent systems.

    \item \textbf{Direct and indirect attacks require different defenses.} Prompt isolation and provenance marking are ineffective against direct injection (where the attacker is the user) but reduce indirect injection ASR by 10--20 percentage points in the LLM experiment. Conversely, tool permission control is effective against both attack types because it operates at the execution layer regardless of where the injection originates.

    \item \textbf{No single defense is sufficient.} The best individual defense (tool permission) still allows 25.0\% of attacks to succeed against the real model. Only the combined configuration achieves the strongest protection across both settings.

    \item \textbf{Security and utility trade off.} Prompt detection introduces 41.0\% utility degradation in the direct setting of the simulation due to false positives on legitimate user requests. In contrast, the same defense causes no degradation in the indirect setting, where attack patterns are more distinctive. This suggests that detection-based defenses are better suited for filtering external content than user input.

    \item \textbf{Data leakage is a distinct risk from attack success.} Even when ASR is relatively low, SDLR can remain high if the agent echoes sensitive data from tool outputs in its response. Prompt isolation and provenance marking are most effective at reducing this specific risk, even when they do not prevent the underlying attack from influencing tool selection.
\end{enumerate}
